\documentclass[12pt]{article}

\usepackage[a4paper,margin=1in]{geometry}
\usepackage{times}
\usepackage[T1]{fontenc}
\usepackage[utf8]{inputenc}
\usepackage{booktabs}
\usepackage{longtable}
\usepackage{array}
\usepackage{enumitem}
\usepackage{titlesec}
\usepackage{xcolor}
\usepackage{hyperref}
\usepackage{setspace}

\onehalfspacing

\titleformat{\section}{\normalfont\Large\bfseries\color{blue!40!black}}{\thesection}{1em}{}
\titleformat{\subsection}{\normalfont\large\bfseries\color{blue!40!black}}{\thesubsection}{1em}{}

\newcolumntype{L}[1]{>{\raggedright\arraybackslash}p{#1}}

\begin{document}

\section*{Chapter 2 --- Requirement Analysis}
\addcontentsline{toc}{section}{Chapter 2 --- Requirement Analysis}

\subsection*{2.1. Stakeholders}

The AR-IMMS system serves several distinct groups of users, each with different needs and objectives:

\begin{longtable}{L{3.2cm} L{10.5cm}}
\toprule
\textbf{Stakeholder} & \textbf{Role / Primary Needs} \\
\midrule
\endhead
Administrator & Manages the overall system: user accounts, role assignment, infrastructure configuration (Site/Room/Rack/Server), and audit log review. \\[4pt]
Operator & Handles day-to-day operations: monitors the dashboard, processes alerts/incidents, approves and closes tickets, oversees overall system status. \\[4pt]
Technician & Performs on-site maintenance: receives tickets, uses AR to scan/identify servers, views telemetry, carries out maintenance work, and updates results. \\[4pt]
Monitoring System & (System actor) --- the Collector and related automated processes that gather and transmit telemetry, and generate alerts when thresholds are exceeded. \\
\bottomrule
\end{longtable}

\textit{Note: this is the core stakeholder list per the outline; you may add ``Management/Executives'' if the report requires a higher-level reporting perspective, or ``Vendor/Equipment Supplier'' if the scope extends to warranty management.}

\subsection*{2.2. System Actors}

This section describes the \textbf{permissions} and \textbf{responsibilities} of each actor identified in Section~2.1 --- forming the basis for the Role-Based Access Control (RBAC) design in Chapter~5.

\paragraph{Administrator}
\begin{itemize}[leftmargin=*]
    \item Permissions: Full system access (CRUD on users, roles, sites/rooms/racks/servers, threshold configuration, audit log viewing).
    \item Responsibilities: Ensuring the infrastructure hierarchy is correctly defined, and managing accounts and permissions securely.
\end{itemize}

\paragraph{Operator}
\begin{itemize}[leftmargin=*]
    \item Permissions: View the dashboard, view/process alerts and incidents, create/assign/close tickets, view reports.
    \item Responsibilities: Continuously monitor system health, coordinate incident response decisions, and approve tickets before closure.
\end{itemize}

\paragraph{Technician}
\begin{itemize}[leftmargin=*]
    \item Permissions: View assigned tickets, use the AR feature to scan devices, view telemetry/alerts for the device being handled, update ticket status, and log maintenance records.
    \item Responsibilities: Perform actual on-site maintenance work and update results accurately and promptly.
\end{itemize}

\paragraph{Monitoring System (Collector)}
\begin{itemize}[leftmargin=*]
    \item Permissions: Send telemetry data to the backend via a dedicated API/channel (typically authenticated with a service token/key rather than a user account).
    \item Responsibilities: Collect data at the correct interval (e.g., every 5 seconds --- see Section~6.5), ensure data completeness, and handle retry/reconnect logic when connectivity is lost (see Section~6.7).
\end{itemize}

\subsubsection*{Permission Matrix --- Module $\times$ Actor}

The following matrix maps each system module to the access level granted to each actor, providing a direct reference for the RBAC design in Chapter~5 (Section~5.5) and for security test-case design in Section~10.6.

\begin{longtable}{L{2.6cm} L{2.6cm} L{2.9cm} L{2.9cm} L{2.6cm}}
\toprule
\textbf{Module} & \textbf{Administrator} & \textbf{Operator} & \textbf{Technician} & \textbf{Monitoring System} \\
\midrule
\endhead
Authentication & Full & Full (own account) & Full (own account) & Token-based (service auth) \\[4pt]
User Management & Full (CRUD) & View only & --- & --- \\[4pt]
Infrastructure Mgmt. & Full (CRUD) & View & View & --- \\[4pt]
Telemetry & View & View & View (assigned device) & Write (submit data) \\[4pt]
Alert & Full (CRUD, config threshold) & View, Acknowledge, Resolve & View (related to own ticket) & Trigger (auto-generated) \\[4pt]
Incident & Full & Create, View, Manage & View (related to own ticket) & --- \\[4pt]
Ticket & Full & Create, Assign, Approve, Close & View (assigned), Update status & --- \\[4pt]
Maintenance & View, Audit & View, Approve & Create, Update (own records) & --- \\[4pt]
Digital Twin & View, Config & View & View (assigned device) & Update status (via telemetry) \\[4pt]
AR (Mobile) & --- & --- & Full (scan, view overlay) & --- \\[4pt]
Audit Log & Full (view all) & View (own actions) & --- & --- \\
\bottomrule
\end{longtable}

\textbf{Legend:} Full = Create/Read/Update/Delete \quad $\cdot$ \quad Manage = create/update/process within role scope, no delete \quad $\cdot$ \quad View = read-only \quad $\cdot$ \quad --- = no access

\textit{Note: The Collector is a system actor, not a human user --- it authenticates via a service token/API key rather than a user session. Technician access is scoped to devices tied to their currently assigned tickets; this scoping rule should be enforced at the API/authorization layer (Section~5.5), not only at the UI level. Each matrix cell is also a candidate test case for Section~10.6 (Security Testing).}

\subsection*{2.3. Functional Requirements}

The following table lists the Functional Requirements (FRs) of the AR-IMMS system, grouped by module. Each item should be assigned a code (FR01, FR02, \ldots) to support the Requirement Traceability Matrix in Section~2.9:

\begin{longtable}{L{2.1cm} L{3.4cm} L{8.2cm}}
\toprule
\textbf{Code} & \textbf{Function Group} & \textbf{Description} \\
\midrule
\endhead
FR01--FR03 & Authentication & Login/logout, JWT-based authentication, password change. \\[4pt]
FR04--FR05 & User Management & Create/edit/delete user accounts, assign roles (Admin/Operator/Technician). \\[4pt]
FR06--FR07 & Infrastructure Management & Manage the Site $\rightarrow$ Room $\rightarrow$ Rack $\rightarrow$ Server $\rightarrow$ Container hierarchy. \\[4pt]
FR08 & Telemetry & Collect and store real-time CPU/RAM/Disk/Network data. \\[4pt]
FR09--FR10 & Alert & Automatically detect anomalies based on thresholds; manage the alert lifecycle (OPEN $\rightarrow$ ACKNOWLEDGED $\rightarrow$ RESOLVED $\rightarrow$ CLOSED). \\[4pt]
FR11 & Incident & Create an incident from an alert when formal handling is required. \\[4pt]
FR12--FR13 & Ticket & Create, assign to a technician, and update ticket status (OPEN $\rightarrow$ IN\_PROGRESS $\rightarrow$ RESOLVED $\rightarrow$ CLOSED). \\[4pt]
FR14 & Maintenance & Record maintenance results (description, time, technician involved). \\[4pt]
FR15 & Digital Twin & Synchronize the physical device's status with its digital representation in real time. \\[4pt]
FR16--FR17 & AR & Scan QR/ArUco markers to identify devices; display an overlay of information (telemetry, alerts, tickets) on the mobile device. \\
\bottomrule
\end{longtable}

\textit{Note: This is a sample table --- cross-check it against your team's actual codebase modules to assign accurate FR codes, since the RTM in Section~2.9 will reference these codes directly.}

\end{document}
